% Cover letter - Quantcast Machine Learning Engineer
\documentclass[letterpaper,11pt]{article}
\usepackage[empty]{fullpage}
\usepackage[hidelinks]{hyperref}
\usepackage[english]{babel}
\usepackage{parskip}
\setlength{\parskip}{6pt}
\usepackage{setspace}
\addtolength{\oddsidemargin}{-0.5in}
\addtolength{\evensidemargin}{-0.5in}
\addtolength{\textwidth}{0.95in}
\addtolength{\topmargin}{-.80in}
\addtolength{\textheight}{1.7in}
\urlstyle{same}
\raggedright
\setstretch{1.0}
\pagestyle{empty}
\begin{document}

Dear Hiring Manager,

I am writing to apply for the Machine Learning Engineer position on the Modeling team at Quantcast. I am completing an MSc in Health Data Science at University College London, where a distinction is expected, after graduating from Imperial College London, and I work as a Data Analyst Intern at TikTok Shop UK. Across my degree and my internships I have spent most of my time on the two things this role asks for: building models, and running the experiments that establish whether they are actually better.

What attracts me most to Quantcast is that the modelling problem is constrained by real engineering limits. Predicting age, gender, viewability, fraud and relevance inside a hundred milliseconds, millions of times a second, is a very dif\kern0pt ferent discipline from f\kern0pt itting a model in a notebook, and the of\kern0pt fer of mentorship from senior scientists to bridge academic concepts with industrial scale is the reason I am applying here rather than somewhere I would be left to my own devices. I enjoy taking an unclear problem, identifying what users actually need, and turning it into a workf\kern0pt low that can be tested and improved.

On probability, statistics and hypothesis testing, my MSc dissertation is the clearest evidence. I compared logistic regression against three tree ensembles for cardiovascular risk prediction on 458,840 UK Biobank participants, judging them on calibration rather than headline accuracy, because a model that ranks well but is badly calibrated cannot support the decision it was built for. I then used doubly robust causal forests to estimate individual-level ef\kern0pt fects, with overlap weighting, sensitivity analysis including E-values, and external validation in a second cohort. I reported null results where I found them.

On running machine learning experiments, at DS Digital Technology I built a Python pipeline calling six LLM APIs with retrieval and search augmentation that produced and scored more than 2,400 responses through an LLM-as-judge framework, against a 246-question benchmark and six-dimension rubric I designed to make quality measurable. I used blinded paired A/B tests to check whether a knowledge graph genuinely improved answer quality rather than assuming it did, and applied NLP and information extraction to more than 4,800 consultation threads.

On the ML ecosystem you list, I use pandas, NumPy, scikit-learn, XGBoost and PyTorch routinely, and Hugging Face Transformers for NLP: I led a UCL group project f\kern0pt ine-tuning ClinicalBERT and PubMedBERT-style models to parse radiology reports into structured features, which outperformed TF-IDF and kNN baselines. At TikTok Shop UK I mine large-scale merchant and transaction data in SQL and Python, and I built an AI agent using Claude Code that automates repetitive steps of the analytics workf\kern0pt low, now used by my team and adjacent teams.

I should be clear about the gaps, because they are relevant to this role specif\kern0pt ically. I have not worked on real-time or high-frequency production systems, I have not written Java, and my exposure to distributed systems, concurrency and low-latency optimisation is academic rather than practical. That engineering side is what I would be learning. My Python, my mathematics and statistics, and my experience designing and evaluating models are well established, and the interest in systems design is genuine rather than stated for the application.

I have the right to work in the UK until January 2029 without sponsorship, and I am based in London.

Thank you for considering my application. I would welcome the opportunity to discuss how my background in machine learning and statistics could support the Modeling team.

\vspace{2pt}
Yours sincerely,

\vspace{8pt}
Zhuoyan Sun

\end{document}
